% ============================================================
%  Chương III – Cấu trúc hạ tầng khóa công khai PKI
%  ~1000 từ · 3 trang
% ============================================================
\chapter{Cấu trúc hạ tầng khóa công khai PKI}
\label{chap:pki-arch}

% Ý chính: PKI không chỉ là kỹ thuật mà là hệ thống tổ hợp
% con người + quy trình + phần mềm + phần cứng + chính sách.


% ----------------------------------------------------------
\section{Các mô hình phân phối khóa công khai}
\label{sec:3.1}
% Ý chính: 4 layer tiến hóa — host-based trust → centralized KDC → Web of Trust → hierarchical PKI
% Mỗi layer giải quyết một vấn đề của layer trước nhưng lại tạo ra vấn đề mới

\subsection{Bài toán phân phối khóa công khai}
% Ý chính: public key cần được công bố rộng rãi nhưng vẫn phải đảm bảo tính bảo mật
% Bốn mục tiêu: xác thực, bảo mật, toàn vẹn, không chối cãi

\subsection{Public Announcement}
% Ý chính: public key được công bố rộng rãi, ai cũng có thể truy cập
% Ví dụ: DNSSEC public key được lưu trong DNS, có thể truy vấn bằng dig
% Vấn đề authentication: làm sao biết public key đó thực sự thuộc về entity nào?

\subsection{Host-based Trust}
% Ý chính: SSH known\_hosts — mỗi host tự quản lý danh sách public key tin cậy
% Vấn đề: O(n²) cặp khóa cho n host; không scale được cho internet

\subsection{Web of Trust (PGP)}
% Ý chính: peer endorsement phi tập trung; mỗi người tự quyết định tin ai
% Thích hợp cộng đồng nhỏ/kỹ thuật; bootstrapping problem; vẫn không scale được \cite{Zimmermann1995}

\subsection{Hierarchical PKI}
% Ý chính: transitive trust từ root xuống leaf; phù hợp Web scale
% Browser trust store (Mozilla NSS, Microsoft, Apple) = danh sách root CA tin cậy
% CA/Browser Forum \cite{CABForum} đặt ra baseline requirements

\subsection{So sánh trade-off}
% Bảng~\ref{tab:trust-models}: centralized vs decentralized
% Hierarchical: dễ audit, dễ revoke, nhưng CA tập trung là single point of failure về trust
% WoT: không SPOF, nhưng trust không đủ transitive và khó dùng cho người phổ thông


% ----------------------------------------------------------
\section{Định nghĩa và mục tiêu PKI}
\label{sec:3.2}
% Ý chính: PKI = tổ hợp phần cứng, phần mềm, con người, chính sách, quy trình
% để quản lý chứng chỉ số trong môi trường mạng mở
% PKI giải quyết bài toán trust bằng cách ràng buộc danh tính với public key thông qua CA
% PKI không chỉ là kỹ thuật mà còn là governance — ai được phép làm CA?


% ----------------------------------------------------------
\section{Các thành phần PKI}
\label{sec:3.3}

\subsubsection*{CA -- Certificate Authority}
% Ý chính: phát hành + thu hồi chứng chỉ; lưu giữ CRL
% Phân tầng: Root CA (offline, air-gapped) → Intermediate CA → Leaf/End-entity
% Quy trình: CSR → xác minh danh tính → ký cert → publish

\subsubsection*{RA -- Registration Authority}
% Ý chính: tách xác minh danh tính (RA) khỏi ký cert (CA)
% Giảm rủi ro, phân quyền trách nhiệm; ví dụ: RA xác minh hộ chiếu trước khi CA ký

\subsubsection*{Repository -- Kho lưu trữ chứng chỉ}
% Ý chính: LDAP directory, HTTP server — nơi client truy vấn cert và CRL public

\subsubsection*{CRL -- Certificate Revocation List}
% Ý chính: danh sách cert đã bị thu hồi trước hạn; reason codes (keyCompromise, cACompromise)
% Hạn chế: CRL có thể rất lớn và stale hàng giờ -- giải pháp: delta CRL \cite{RFC5280}

\subsubsection*{OCSP -- Online Certificate Status Protocol}
% Ý chính: kiểm tra realtime trạng thái từng cert theo serial number \cite{RFC6960}
% OCSP Stapling \cite{RFC6066}: server pre-fetch response, giải quyết privacy và giảm latency
% Bảng~\ref{tab:revocation}: so sánh CRL vs OCSP vs OCSP Stapling


% ----------------------------------------------------------
\section{Chu kỳ sống của chứng chỉ}
\label{sec:3.4}
% Ý chính: Tạo keypair → CSR → xác minh danh tính → phát hành → sử dụng → gia hạn / thu hồi
% Certificate lifecycle automation (ACME, SCEP, EST) ngày càng quan trọng
% Hình~\ref{fig:cert-lifecycle}: sơ đồ lifecycle

